\subsection{Introduction}
Consider a finite collection of stochastic alternatives with unknown mean rewards. A similarity graph can express prior information that two alternatives should perform similarly. An observation at one alternative can then inform estimates at another without directly observing its reward. The similarity selection index introduced by Sun, Li, and Fu~\cite{collected14}, and developed by Zhou, Fu, and Ryzhov~\cite{collected15}, formalizes this intuition through Laplacian smoothing. Their objective is final selection. The present question is how such information changes the cumulative regret of sequential decisions.

For classical stationary stochastic bandits, the principal targets are established. Worst-case expected regret for bounded rewards has optimal order $\sqrt{KT}$, and fixed-instance asymptotic lower bounds admit matching algorithms for important distribution families~\cite{collected20,collected21,collected22}. Consequently, a graph-based algorithm must be assessed against a restricted class of reward vectors that encodes what the graph tells the learner. A lower regret bound is possible because fewer alternative environments remain plausible, rather than because smoothing creates additional observations.

The phrase ``better alignment'' needs care. A graph can preserve the ordering of all means while providing no improvement in the information needed to distinguish them. Conversely, a graph may accurately identify several groups of nearly equivalent suboptimal alternatives, allowing a whole group to be rejected after a small pooled sample. Its usefulness depends on its geometry, the amount of reward variation allowed within it, and the gaps between groups. No scalar measure of visual resemblance or ranking agreement guarantees monotone regret improvement for arbitrary algorithms and graphs.

Our principal result shows why the full graph geometry matters. A complete graph with edge weights $1/m$ and a path with edge weights $(m-1)/2$ both have resistance diameter two. At homogeneous suboptimal means and supplied energy radius equal to the gap, every clique arm can independently become optimal within the model class. The path excludes these independent deviations: an average of its endpoints controls every interior mean. The clique requires an information coefficient proportional to $m$, while a concrete endpoint-based policy on the path has no linear dependence on $m$. This is a separation between graph-energy classes with identical diameter summaries, not a claim that spectral bandits or graph-aware exploration are new.

We develop the result in three steps. First, we express the information constant as an experimental design over least informative graph-energy alternatives. A finite-strength resolvent evaluates its inner optimization. Second, we show that the first-order cost of loosening alignment is governed by the minimum enclosing radius in the resistance embedding, equivalently an established maximum graph variance quantity. Third, a capped graph-design elimination algorithm uses this geometry and falls back to ordinary arm exploration on competitive components. We prove its regret bound, the achieved path separation, and a quantitative comparison under certified Laplacian perturbations. A simpler diameter-based pooling policy remains as a baseline and clarifies the connection to the original selection index.

The distinctions among results remain explicit. The information design gives the exact value of a lower-bound program for homogeneous components. The radius expansion describes that program's local sensitivity. The proposed policy achieves the path's improved arm-count dependence with conservative constants; it is not proved to attain the information constant for positive uncertainty. A separate standard lower-bound argument explains why unrestricted graph validation cannot preserve the reduced logarithmic constant for free.

\subsection{Related research and contribution boundary}
The original spectral selection index minimizes a squared estimation term plus a graph-Laplacian penalty~\cite{collected14}. Subsequent work develops sequential allocation for selecting the best alternative and permits updating similarity information~\cite{collected15}. Those allocation objectives differ from cumulative reward: a positive limiting sampling fraction assigned to a suboptimal arm produces linear cumulative regret. We retain the graph estimator's interpretation while changing the allocation objective.

Graph-smooth cumulative-regret bandits already appear in SpectralUCB and related algorithms~\cite{collected1,collected17}. They use graph spectral structure to reduce effective statistical complexity. We therefore do not claim that graph arms, Laplacian regularization, or a spectral confidence index constitute a new problem or algorithmic principle.

Clustered bandits are also established. Carlsson, Dubhashi, and Johansson~\cite{collected18} analyze hierarchical Thompson sampling and show how clustering quality affects regret under dominance assumptions. Gore and Chaporkar's Clus-UCB preprint~\cite{collected19} directly studies known clusters with bounded mean differences, derives an information lower bound, and proposes an index that shares observations within clusters. This is a close comparator for our pooling analysis. The preprint is cited as such; no publication status is assumed.

Our lower-bound starting point is the structured-bandit change-of-measure program of Combes, Magureanu, and Proutiere~\cite{collected11}. The maximum graph variance construction underlying our resistance design is established by Devriendt, Martin-Gutierrez, and Lambiotte~\cite{collected24}. Graph-aware pure exploration and resistance-based influence measures also appear in GRUB~\cite{collected2}. We do not claim to introduce experimental design, resistance geometry, or eliminating unsampled graph arms. The candidate contributions are the graph-energy information sensitivity, its explicit separation from diameter compression, and the accompanying cumulative-regret analysis and certified perturbation comparison.

The clique formula needs particularly careful positioning. For homogeneous clusters, Clus-UCB's information expression has the form $m\Gamma/[d+(m-1)b]$. Substituting Gaussian divergences gives the same algebraic shape as our clique constant with effective width $S/\sqrt{w(m-1)}$. This is our comparison, not a Gaussian theorem proved in that Bernoulli paper. It motivates making the general-graph separation, rather than the clique interpolation, the central result. We inspected Clus-UCB v2, including its discussion of incorrect widths~\cite{collected19}.

Model misspecification and unknown error are established topics in kernelized bandits~\cite{collected25}. Gap-adjusted approximation error~\cite{collected26} and graph-specific residual and eigenspace perturbation bounds~\cite{collected27} further constrain broad novelty claims. Our perturbation result assumes a certified comparison to a reference graph. We do not claim to learn that comparison from unrestricted selected-arm feedback. Independent mathematical review and a broader literature comparison remain necessary before asserting publication priority.

Similarity graphs must also be distinguished from feedback graphs. In our setting, a pull reveals only one reward. An edge restricts the possible means; it does not supply observations of adjacent arms. Observation noise is independent in the experiments. Dependence of the unknown means through a model class is sufficient for information sharing.

\subsection{Model and definitions of alignment}
\subsubsection{Bandit observations}
There are $K$ arms. At time $t$, an adapted policy selects $A_t$ and observes
\begin{equation}
 Y_t=\mu_{A_t}+\xi_t,\qquad \mu\in[0,1]^K.
 \label{alignment_paper:eq:model}
\end{equation}
Conditional on the past and the chosen arm, the noise has zero mean and satisfies
$\E[\exp(u\xi_t)\mid\mathcal F_{t-1},A_t]\le\exp(u^2\sigma^2/2)$ for every $u\in\R$.
Independent iid arm streams with common known sub-Gaussian proxy $\sigma$ suffice. Our lower bounds specialize to independent $N(\mu_i,\sigma^2)$ rewards with known variance. Define
\begin{equation}
 \mu^\star=\max_i\mu_i,\quad \Delta_i=\mu^\star-\mu_i,
 \quad \mathcal R_T=\sum_{t=1}^T\Delta_{A_t},\quad
 R_T=\E[\mathcal R_T].
\end{equation}
All regret is measured with the original means, even when the selection index is smoothed.

\subsubsection{Ordinal alignment}
Let $W$ be symmetric with nonnegative off-diagonal weights, write $\one$ for the all-ones vector, and let
$L=\diag(W\one)-W$. Write
\begin{equation}
 Q_\lambda=(I+\lambda L)^{-1},\qquad z_\lambda=Q_\lambda\bar Y.
 \label{alignment_paper:eq:sindex}
\end{equation}
We call a filter ordinally aligned at $\mu$ when $Q_\lambda\mu$ preserves its strict reward comparisons. This describes the order-preservation property relevant here. It does not replace the specific structural definitions of an aligned graph used in the earlier selection literature. In particular, the original WSC definition includes equal-degree and reward-distance conditions, while the later dissertation relaxes the structural condition~\cite{collected14,collected16}.

\begin{proposition}[Order preservation need not improve a comparison]
For the complete graph with common edge weight $w>0$,
\begin{equation}
 Q_\lambda=qI+(1-q)\frac{\one\one^\top}{K},\qquad
 q=(1+\lambda wK)^{-1}.
\end{equation}
For every vector $y$, $z_i-z_j=q(y_i-y_j)$. Thus the smoothed and raw sample means select exactly the same maximizer. With independent Gaussian sample means based on equal fixed counts $n$, both the true pairwise gap and its standard deviation are multiplied by $q$, leaving the signal-to-noise ratio unchanged.
\label{alignment_paper:prop:ordinal}
\end{proposition}
\begin{proof}
The complete-graph Laplacian is $w(KI-\one\one^\top)$. Its eigenvalue is zero on constants and $wK$ on their orthogonal complement. Inverting on these two subspaces gives the displayed expression. Pairwise differences remove the common-average term. Their raw variance is $2\sigma^2/n$, and their filtered variance is $q^2 2\sigma^2/n$.
\end{proof}
This example shows the limitation of rank preservation as a performance claim. It does not claim that every algorithm using a complete graph behaves identically to UCB.

\subsubsection{Quantitative alignment through graph energy}
Choose a fixed partition $\mathcal P=\{C_1,\ldots,C_M\}$ of the arms. A graph may already have these connected components, or a partition may be constructed and cross-partition edges removed. Let $L_c$ denote the resulting connected component Laplacian; singletons are permitted. Removing edges is part of the model construction, not an assumption that a connected original graph has zero cross edges.

The learner is supplied with valid local certificates
\begin{equation}
 \mu_c^\top L_c\mu_c\le S_c^2.
 \label{alignment_paper:eq:energy}
\end{equation}
The corresponding model class is
\begin{equation}
 \mathcal M_{\mathcal P,L,S}=\{\nu\in[0,1]^K:
       \nu_c^\top L_c\nu_c\le S_c^2\ \text{for all }c\}.
 \label{alignment_paper:eq:class}
\end{equation}
The certificates are prior information available to the policy. Their true values cannot silently be computed from unknown rewards during a real deployment.

For a connected component, define effective resistance
$r_{ij}^{\rm eff}=(e_i-e_j)^\top L_c^\dagger(e_i-e_j)$ and resistance diameter
$D_c^{\rm eff}=\max_{i,j\in C_c}r_{ij}^{\rm eff}$. Use $D_c^{\rm eff}=0$ for a singleton.

\begin{lemma}[Energy certifies reward diameter]
If \eqref{alignment_paper:eq:energy} holds, then
\begin{equation}
 |\mu_i-\mu_j|\le S_c\sqrt{r_{ij}^{\rm eff}},\qquad
 \max_{i\in C_c}\mu_i-\min_{i\in C_c}\mu_i
 \le\varepsilon_c:=\min\{1,S_c\sqrt{D_c^{\rm eff}}\}.
 \label{alignment_paper:eq:width}
\end{equation}
\label{alignment_paper:lem:resistance}
\end{lemma}
\begin{proof}
The vector $e_i-e_j$ is orthogonal to constants. On this subspace, factor its inner product with $\mu_c$ through $L_c^{\dagger/2}$ and $L_c^{1/2}$. Cauchy--Schwarz gives
$|(e_i-e_j)^\top\mu_c|^2\le
((e_i-e_j)^\top L_c^\dagger(e_i-e_j))(\mu_c^\top L_c\mu_c)$.
Bounded means additionally give diameter at most one.
\end{proof}
The product $S_c\sqrt{D_c^{\rm eff}}$ is invariant to a common rescaling of all component weights, provided the energy certificate is rescaled consistently. Shrinking weights without changing the information they represent is therefore not a substantive improvement in alignment.

Let $v_c=\max_{i\in C_c}\mu_i$, $\Gamma_c=\mu^\star-v_c$, and
$\mathcal O=\{c:\Gamma_c=0\}$. The relevant dimensionless uncertainty for rejecting a suboptimal component is $\varepsilon_c/\Gamma_c$. Neither $\Gamma_c$ nor this ratio is an algorithm input.

\subsection{From the spectral index to a pooling policy}
\subsubsection{Count-weighted spectral pooling}
Let $n_i$ and $\bar Y_i$ denote observed counts and sample means. The count-weighted analogue of the selection index solves
\begin{equation}
 \widehat f_\lambda\in\argmin_f
       \sum_{i:n_i>0}n_i(f_i-\bar Y_i)^2+\lambda f^\top Lf,
 \label{alignment_paper:eq:weighted}
\end{equation}
where $L$ is the block-diagonal Laplacian of the chosen partition. After every component has an observation, the solution is unique: a vector in the null spaces of both $\diag(n_i)$ and $L$ must vanish on every component. As $\lambda\to\infty$, nonconstant component modes are suppressed. Minimizing the remaining squared-error objective over constants gives
\begin{equation}
 \widehat f_{\infty,i}=\widehat p_c
 =\frac{\sum_{s:A_s\in C_c}Y_s}{N_c},\qquad i\in C_c,
 \quad N_c=\sum_{i\in C_c}n_i.
 \label{alignment_paper:eq:pool}
\end{equation}
This is a spectral pooling limit. It weights observations by counts; it is not the unweighted average of available arm averages. For unequal counts it differs from applying a fixed $Q_\lambda$ to the raw means. We do not invoke the earlier order-preservation theorem for this modified estimator.

The conditional mean of the pooled observation sequence is the average of the actual means of the sampled arms. Because selection is adaptive, $\widehat p_c$ is not an iid estimate of a fixed component mean when its arms differ. The component diameter certificate controls this drift.

\subsubsection{Algorithm}
Assume horizon $T\ge M$ and confidence level $\delta\in(0,1)$. Put
\begin{equation}
 \ell=\log\frac{2(K+M)T}{\delta},\qquad
 b(n)=\sigma\sqrt{\frac{2\ell}{n}},\quad b(0)=+\infty.
 \label{alignment_paper:eq:radius}
\end{equation}
For an unsampled arm set $u_i=+\infty$; otherwise set
$u_i=\bar Y_i+b(n_i)$. Define
\begin{equation}
 i_c\in\argmax_{i\in C_c}u_i,\qquad
 U_c=\min\{u_{i_c},\widehat p_c+b(N_c)+\varepsilon_c\}.
 \label{alignment_paper:eq:index}
\end{equation}
The policy, called spectral pooling UCB (SP-UCB) in this draft, is:
\begin{enumerate}
 \item Pull one predetermined representative of each component.
 \item At subsequent rounds select $c_t\in\argmax_c U_c$.
 \item Pull $A_t=i_{c_t}$ and update its arm and component statistics.
\end{enumerate}
Tie breaking uses a fixed arm ordering. There is no all-arm initialization requirement: a clearly suboptimal component can be rejected before all its members are sampled. An uncertain optimal component retains ordinary arm exploration through $i_c$.

The minimum in \eqref{alignment_paper:eq:index} combines two valid bounds on the same target $v_c$. Its validity depends on the supplied certificate. Taking a minimum with an ordinary UCB is not protection against a false graph assertion.

\subsubsection{Confidence under adaptive pooling}
For the $n$th observation from component $c$, let $B_{c,n}$ be the average of the means of the $n$ arms actually pulled there, including repetitions. A uniform event $\mathcal E$ is
\begin{equation}
 |\bar Y_i-\mu_i|\le b(n_i),\qquad
 |\widehat p_c-B_{c,N_c}|\le b(N_c)
 \label{alignment_paper:eq:event}
\end{equation}
for all nonzero arm and component counts through $T$.

\begin{lemma}[Uniform confidence]
Under \eqref{alignment_paper:eq:model}, $\Prob(\mathcal E)\ge1-\delta$. On this event, every $U_c$ in \eqref{alignment_paper:eq:index} is an upper bound on $v_c$.
\label{alignment_paper:lem:confidence}
\end{lemma}
\begin{proof}
For a fixed arm or component, the selected-noise sum $Z_t$ and count $N_t$ give the exponential supermartingale
$\exp(uZ_t-u^2\sigma^2N_t/2)$. Stop it at the $n$th corresponding pull or at $T$. On the event that the $n$th pull occurs, the upper-tail probability for $Z\ge nb(n)$ is at most $e^{-\ell}$, using $u=b(n)/\sigma^2$. Apply the same argument to the lower tail. A union bound over $K+M$ processes and $n=1,\ldots,T$ gives failure probability at most $2(K+M)Te^{-\ell}=\delta$. This argument does not treat a stopped sample average as Gaussian.

For a component, $v_c-\varepsilon_c\le B_{c,N_c}\le v_c$. Thus $\widehat p_c+b(N_c)+\varepsilon_c\ge v_c$ on $\mathcal E$. The largest ordinary arm upper bound also exceeds $v_c$. Their minimum remains an upper bound.
\end{proof}

\subsubsection{Finite-time regret}
Define the ordinary component bound
\begin{equation}
 A_c=\sum_{i\in C_c:\Delta_i>0}
        \left(\Delta_i+\frac{8\sigma^2\ell}{\Delta_i}\right).
 \label{alignment_paper:eq:Ac}
\end{equation}
For suboptimal components define
\begin{equation}
 H_c=\begin{cases}
 (\Gamma_c+\varepsilon_c)
 \left(1+\dfrac{8\sigma^2\ell}{(\Gamma_c-\varepsilon_c)^2}\right),
       &\varepsilon_c<\Gamma_c,\\
 +\infty,&\varepsilon_c\ge\Gamma_c.
 \end{cases}
 \label{alignment_paper:eq:Hc}
\end{equation}

\begin{theorem}[Alignment-sensitive finite-time guarantee]
If the graph-derived or directly supplied diameter certificates are valid, SP-UCB satisfies, on $\mathcal E$,
\begin{equation}
 \mathcal R_T\le
     \sum_{c\in\mathcal O}A_c+
     \sum_{c\notin\mathcal O}\min\{A_c,H_c\}.
 \label{alignment_paper:eq:regret}
\end{equation}
Consequently, expected regret is at most the right side plus $\delta T$.
\label{alignment_paper:thm:upper}
\end{theorem}
\begin{proof}
At any post-initialization round selecting arm $i$ in component $c$, optimism of an optimal component gives
$\mu^\star\le U_c\le u_i\le\mu_i+2b(n_i)$ whenever $n_i>0$.
Thus a suboptimal arm can be selected with a positive pre-pull count only if
$n_i\le8\sigma^2\ell/\Delta_i^2$.
Its first pull, including a possible initialization pull, contributes the extra one. Therefore
$n_i(T)\le1+8\sigma^2\ell/\Delta_i^2$, yielding $A_c$.

The pooled index also gives
$\mu^\star\le U_c\le v_c+2b(N_c)+\varepsilon_c$.
If $\Gamma_c>\varepsilon_c$, a component can be selected only with
$N_c\le8\sigma^2\ell/(\Gamma_c-\varepsilon_c)^2$ beforehand, except for its first initialization pull. Hence
$N_c(T)\le1+8\sigma^2\ell/(\Gamma_c-\varepsilon_c)^2$.
Each such pull has regret at most $\Gamma_c+\varepsilon_c$, giving $H_c$.
Both bounds hold simultaneously, so take their minimum. On the failure event, bounded means give $\mathcal R_T\le T$.
\end{proof}

\begin{corollary}[Improved certificates tighten the component guarantee]
For a fixed partition and reward vector, $H_c$ is increasing in $\varepsilon_c$ on $[0,\Gamma_c)$. Thus tightening a valid certificate cannot worsen the bound in \eqref{alignment_paper:eq:regret}. If $\varepsilon_c=0$, every arm in that component has gap $\Gamma_c$, and the leading logarithmic terms of $A_c$ and $H_c$ differ by a factor $|C_c|$.
\end{corollary}
\begin{proof}
The derivative of $(\Gamma+\varepsilon)/(\Gamma-\varepsilon)^2$ is
$(3\Gamma+\varepsilon)/(\Gamma-\varepsilon)^3>0$; the nonlogarithmic term is increasing as well. At zero width, $A_c=|C_c|\Gamma_c+8\sigma^2\ell|C_c|/\Gamma_c$, whereas $H_c=\Gamma_c+8\sigma^2\ell/\Gamma_c$.
\end{proof}
This is monotonicity of a proved upper bound, not a pathwise ordering between two policies with different certificates. Changing the partition also changes component gaps and the number of components; no universal monotonicity assertion follows from reducing an empirical graph-error score.

\subsection{How alignment changes the information lower bound}
\subsubsection{General graph classes}
For this section assume Gaussian rewards with known common variance $\sigma^2$, a unique optimal arm $i^\star$, and $\mu^\star<1$. An algorithm is uniformly efficient on a class if its regret is $o(T^a)$ for every $a>0$ at every member of that class. Define confusing alternatives
\begin{equation}
 \mathcal B(\mu)=\{\nu\in\mathcal M_{\mathcal P,L,S}:
       \nu_{i^\star}=\mu^\star,\ \max_j\nu_j>\mu^\star\}.
\end{equation}
The established structured-bandit information principle~\cite{collected11} yields the program
\begin{align}
 C_{L,S}(\mu)=\inf_{x_i\ge0}\quad&\sum_{i\ne i^\star}x_i\Delta_i,\label{alignment_paper:eq:info}\\
 \text{subject to}\quad&
 \sum_{i\ne i^\star}x_i\frac{(\mu_i-\nu_i)^2}{2\sigma^2}\ge1
       \quad\text{for every }\nu\in\mathcal B(\mu).\nonumber
\end{align}
In particular,
\begin{equation}
 \liminf_{T\to\infty}\frac{R_T}{\log T}\ge C_{L,S}(\mu).
 \label{alignment_paper:eq:lower}
\end{equation}
The normalization $x_i$ represents the number of samples of arm $i$ per unit $\log T$. Alternatives agreeing on the optimal arm cannot be rejected by cost-free optimal-arm observations alone.

The following monotonicity is a property of the information available, independent of a particular algorithm. If two supplied classes contain the true vector and
$\mathcal M_1\subseteq\mathcal M_2$, then $C_{\mathcal M_1}(\mu)\le C_{\mathcal M_2}(\mu)$. The first class has fewer alternative constraints in \eqref{alignment_paper:eq:info}. Decreasing valid $S_c$ with $L_c$ fixed gives this nesting. Arbitrarily changing graph edges does not necessarily do so.

For completeness, the change-of-measure reasoning is as follows. Comparing $\mu$ with a confusing $\nu$, the KL divergence of histories is
$\sum_i\E_\mu[n_i(T)](\mu_i-\nu_i)^2/(2\sigma^2)$.
Consider the event that the original best arm is played at least $T/2$ times. Uniform efficiency makes its probability approach one under $\mu$ and at most $o(T^{-1+a})$ under $\nu$, for every $a>0$. The binary KL divergence of this event is consequently at least $(1-o(1))\log T$. This gives each information constraint; normalizing sample counts and passing to a finite convergent subsequence of any bounded-cost allocation gives \eqref{alignment_paper:eq:lower}. If normalized regret is unbounded, the inequality is immediate.


\subsubsection{A graph-specific information design}
The diameter bound \eqref{alignment_paper:eq:width} can discard statistically useful geometry. Consider a homogeneous suboptimal component with connected Laplacian $L$, $m\ge2$ arms of true mean $a$, gap $\Gamma=\mu^\star-a$, and supplied energy budget $S^2$. Let $s=S/\Gamma$ and let $\mathscr P_m=\{p\ge0:\one^\top p=1\}$. Define
\begin{align}
 F_L(p,s)&=\min_{i\in\{1,\ldots,m\}}\ \min_{z\in\mathcal Z_i(s)}
                 \sum_jp_jz_j^2,\label{alignment_paper:eq:designF}\\
 \mathcal Z_i(s)&=\{z\in[0,1]^m:z_i=1,\ z^\top Lz\le s^2\},\qquad
 F_L^\star(s)=\max_{p\in\mathscr P_m}F_L(p,s).\nonumber
\end{align}
Here $p$ allocates observations, rather than graph edge weights. The least informative alternative $z$ describes a displacement from the true constant vector, normalized by the gap.

\begin{theorem}[Exact information design for a homogeneous graph component]
When the optimal component is a singleton and all other true components are homogeneous, the program \eqref{alignment_paper:eq:info} separates into component contributions
\begin{equation}
 C_{L,S,c}(\mu)=\frac{2\sigma^2}{\Gamma_c F_{L_c}^\star(S_c/\Gamma_c)}.
 \label{alignment_paper:eq:generaldesignconstant}
\end{equation}
For a fixed graph, $F_L^\star$ is nonincreasing in $s$, lies in $[1/m,1]$, and equals one at $s=0$. For a fixed $s$, $F_L(p,s)$ is concave and continuous in $p$.
\label{alignment_paper:thm:generaldesign}
\end{theorem}
\begin{proof}
Write a component allocation as $x=Xp$, where $X=\sum_jx_j$. Clipping any alternative's displacements to $[0,\Gamma]$ decreases both their squared information cost and graph energy, because clipping is a contraction on every edge. Scaling a displacement that exceeds $\Gamma$ gives a boundary alternative with one coordinate equal to $\Gamma$. Conversely, adding a common arbitrarily small positive displacement to a boundary alternative makes it strictly confusing without changing energy; $\mu^\star<1$ keeps it inside the mean domain. The infimal information rate is therefore $X\Gamma^2F_L(p,s)/(2\sigma^2)$. Minimizing the regret cost $X\Gamma$ proves \eqref{alignment_paper:eq:generaldesignconstant}. Changes in other components cannot make this alternative harder to distinguish, so component costs add.

The sets $\mathcal Z_i(s)$ are compact and nonempty, containing $\one$. Each candidate objective is linear in $p$, which gives concavity and continuity of their minimum. Increasing $s$ expands the candidate sets. For uniform $p$, $z_i=1$ implies objective at least $1/m$, while $\one$ gives the upper bound one for every $p$. At zero energy, connectedness forces $z=\one$.
\end{proof}
This theorem evaluates the established information framework~\cite{collected11}; it does not claim that an arbitrary maximizer of \eqref{alignment_paper:eq:designF} immediately gives an implementable asymptotically optimal bandit policy.

\begin{lemma}[Resolvent oracle for the least informative alternative]
For $p_j>0$ and $s>0$, with $D_p=\diag(p)$, the inner minimum in \eqref{alignment_paper:eq:designF} equals
\begin{equation}
 \sup_{\tau\ge0}\left\{
 \frac{1}{[(D_p+\tau L)^{-1}]_{ii}}-\tau s^2\right\}.
 \label{alignment_paper:eq:resolventoracle}
\end{equation}
For a fixed multiplier, its minimizing vector is
$z=(D_p+\tau L)^{-1}e_i/[(D_p+\tau L)^{-1}]_{ii}$.
\label{alignment_paper:lem:resolvent}
\end{lemma}
\begin{proof}
Relax the box constraint temporarily and dualize the energy inequality. Minimizing $z^\top(D_p+\tau L)z$ subject to $z_i=1$ gives the displayed inverse formula by a Lagrange multiplier. The free-coordinate equations and the maximum principle imply $0\le z_j\le1$, so removing the box constraint does not alter this minimum. The constant vector is strictly feasible for the energy inequality when $s>0$, giving strong convex duality. For $s=0$ the value follows by the limit $\tau\to\infty$. Designs with zero entries are handled by limits of positive designs, using continuity in $p$.
\end{proof}
With uniform design, $(D_p+\tau L)^{-1}=m(I+m\tau L)^{-1}$. Thus the finite-strength spectral resolvent used in the selection index also appears in an exact information calculation, not just in the infinite-smoothing pooling limit. This is a variational connection, not a regret guarantee for selecting the largest fixed-strength index.

\subsubsection{Resistance geometry determines local alignment sensitivity}
Let $\mathcal D$ be the matrix of pairwise effective resistances. For a design $p$, define
\begin{equation}
 \kappa_L(p)=\max_i\sqrt{(e_i-p)^\top L^\dagger(e_i-p)},\qquad
 \rho_L=\min_{p\in\mathscr P_m}\kappa_L(p).
 \label{alignment_paper:eq:radiusdesign}
\end{equation}
The vectors $v_i=L^{\dagger/2}e_i$ embed the arms into Euclidean space with squared distances $\|v_i-v_j\|^2=\mathcal D_{ij}$. The quantity $\rho_L$ is the radius of their smallest enclosing ball. Its center lies in their convex hull, so it has a probability-vector representation. This geometry and the equivalent maximum graph variance problem are established~\cite{collected24}; the next theorem uses them to quantify bandit information sensitivity.

\begin{lemma}[Computable resistance design and a stopping certificate]
\begin{equation}
 \rho_L^2=\max_{p\in\mathscr P_m}\frac12p^\top\mathcal Dp.
 \label{alignment_paper:eq:maxvariance}
\end{equation}
For every $p$, $V(p)=p^\top\mathcal Dp/2$ is a lower bound on $\rho_L^2$, and $\kappa_L(p)^2$ is an upper bound, with gap
\begin{equation}
 \kappa_L(p)^2-V(p)=\max_i(\mathcal Dp)_i-p^\top\mathcal Dp.
 \label{alignment_paper:eq:designgap}
\end{equation}
\label{alignment_paper:lem:designradius}
\end{lemma}
\begin{proof}
For any center $v(p)$, the squared distance to arm $i$ is
$(\mathcal Dp)_i-V(p)$. The variance of any distribution $q$ about its own mean is $V(q)$; its average squared distance about any other center is $V(q)+\|v(q)-v(p)\|^2$. Consequently every enclosing ball has squared radius at least $\max_qV(q)$. At a maximum-variance distribution $p^\star$, simplex optimality gives $(\mathcal Dp^\star)_i\le(p^\star)^\top\mathcal Dp^\star$ for every $i$, with equality on its support. Its center therefore has maximum squared distance $V(p^\star)$, proving equality. Substituting the distance formula gives the gap certificate.
\end{proof}
The radius is invariant in the relevant units: rescaling $L$ by $b>0$ changes $\rho_L$ by $b^{-1/2}$ and a consistently scaled $S$ by $b^{1/2}$. A numerical design need not achieve exact optimality; its computed upper radius gives a valid, possibly looser, certificate.

\begin{theorem}[First-order graph alignment sensitivity]
Let $D=\max_{i,j}\mathcal D_{ij}$. For $s\sqrt D\le1$,
\begin{equation}
 1-2s\rho_L\ \le F_L^\star(s)
 \le 1-2s\rho_L+s^2D.
 \label{alignment_paper:eq:firstorderF}
\end{equation}
For every $s\ge0$,
\begin{equation}
 F_L^\star(s)\ge(1-s\rho_L)_+^2.
 \label{alignment_paper:eq:jensenF}
\end{equation}
In particular, as $S/\Gamma\downarrow0$,
\begin{equation}
 \frac{C_{L,S,c}(\mu)}{2\sigma^2/\Gamma}
 =1+2\frac S\Gamma\rho_L
   +O\!\left(\frac{S^2}{\Gamma^2}D\right).
 \label{alignment_paper:eq:sensitivity}
\end{equation}
\label{alignment_paper:thm:sensitivity}
\end{theorem}
\begin{proof}
For $z\in\mathcal Z_i(s)$, energy Cauchy--Schwarz gives
$1-p^\top z\le s\sqrt{(e_i-p)^\top L^\dagger(e_i-p)}$.
Jensen's inequality and nonnegativity then give
$\sum_jp_jz_j^2\ge(1-s\kappa_L(p))_+^2$. Optimize over $p$ for \eqref{alignment_paper:eq:jensenF}.

For the sharper expansion, write $z=\one-su$ with $u_i=0$ and $u^\top Lu\le1$. The largest possible $p^\top u$ is $\kappa_i(p)=\sqrt{(e_i-p)^\top L^\dagger(e_i-p)}$. A maximizer is obtained from the inverse grounded Laplacian applied to $p_{-i}$ and normalized to unit energy. Its coordinates are nonnegative by the maximum principle. Energy Cauchy--Schwarz additionally gives $u_j\le\sqrt{\mathcal D_{ij}}\le\sqrt D$, so $z$ is box-feasible when $s\sqrt D\le1$. The objective equals $1-2sp^\top u+s^2\sum_jp_ju_j^2$. Its nonnegative final term gives the lower bound $1-2s\kappa_L(p)$; taking a maximizing $u$ for an arm attaining $\kappa_L(p)$ gives the upper bound $1-2s\kappa_L(p)+s^2D$. Optimize over $p$. The information expansion follows by inversion for sufficiently small $s$, using $\rho_L^2\le D$.
\end{proof}
The leading sensitivity is therefore a graph-wide resistance radius, not the resistance diameter or the number of arms alone. A design minimizing that radius is first-order optimal for this information calculation. No first-order optimality claim is made for the regret constant of the algorithm introduced below.

\subsubsection{Certified graph perturbations}
\begin{proposition}[An explicit price for a certified Laplacian error]
Suppose, on the subspace orthogonal to constants,
\begin{equation}
 (1-\eta)L\preceq\widehat L\preceq(1+\eta)L,\qquad 0\le\eta<1.
 \label{alignment_paper:eq:lapperturb}
\end{equation}
If $S$ is valid for $L$, then $\widehat S=\sqrt{1+\eta}\,S$ is valid for $\widehat L$, and
\begin{equation}
 \widehat S\rho_{\widehat L}
 \le S\rho_L\sqrt{\frac{1+\eta}{1-\eta}}.
 \label{alignment_paper:eq:perturbradius}
\end{equation}
The corresponding homogeneous-component information constants obey
\begin{equation}
 C_{L,S,c}\le C_{\widehat L,\widehat S,c}
 \le C_{L,\,S\sqrt{(1+\eta)/(1-\eta)},c}.
 \label{alignment_paper:eq:perturbinfo}
\end{equation}
\label{alignment_paper:prop:perturb}
\end{proposition}
\begin{proof}
The upper Loewner inequality certifies the true energy. Inverting the inequalities on the constant-free subspace bounds every resistance quadratic form between factors $(1+\eta)^{-1}$ and $(1-\eta)^{-1}$ of the original form. Minimize their maximum over designs to obtain \eqref{alignment_paper:eq:perturbradius}. The original energy ball of radius $S$ is contained in the estimated-graph energy ball of radius $\widehat S$, which in turn is contained in the original ball of radius $S\sqrt{(1+\eta)/(1-\eta)}$. Apply information monotonicity.
\end{proof}
For Laplacians with the same constant null space, an operator error bound
$\|\widehat L-L\|_{\rm op}\le\eta\lambda_2(L)$ suffices for \eqref{alignment_paper:eq:lapperturb}. This connects a certified graph error to both an algorithmic bias radius and an information-class comparison. The true reference graph and error bound are additional information; they are not assumed obtainable for free from bandit observations. Related spectral perturbation and mismatch questions already arise in graph-contextual bandits~\cite{collected27}.

\subsubsection{Graphs with identical diameters can have different exploration costs}
Compare two components with the same number $m\ge3$ of arms:
\begin{enumerate}
 \item A complete graph with edge weight $1/m$ has $\mathcal D_{ij}=2$ for distinct arms and $\rho_L^2=(m-1)/m$, attained by the uniform design.
 \item A path with edge weight $(m-1)/2$ has $\mathcal D_{ij}=2|i-j|/(m-1)$ and $\rho_L^2=1/2$, attained by placing half the design mass at each endpoint.
\end{enumerate}
For the path, an embedding with orthogonal successive increments has all its vertices at squared distance $1/2$ from the endpoint midpoint. The endpoints are squared distance two apart, so no smaller ball can enclose them. For the clique, the regular simplex has radius squared $(m-1)/m$. Both graphs have resistance diameter two.

\begin{corollary}[Diameter compression loses an unbounded information factor]
At a homogeneous instance with supplied $S=\Gamma$, the clique contribution is the unstructured value
\begin{equation}
 C_{\rm clique}=\frac{2\sigma^2m}{\Gamma},
 \label{alignment_paper:eq:cliqueseparation}
\end{equation}
whereas the path contribution satisfies
\begin{equation}
 C_{\rm path}\le
 \frac{2\sigma^2}{\Gamma(1-1/\sqrt2)^2}.
 \label{alignment_paper:eq:pathinformation}
\end{equation}
Their ratio is at least $m(1-1/\sqrt2)^2$, which grows without bound with $m$. The number of arms, true means, supplied energy budget, and resistance diameter are identical.
\label{alignment_paper:cor:informationseparation}
\end{corollary}
\begin{proof}
In the clique, a single-coordinate displacement of $\Gamma$ costs energy $\Gamma^2(m-1)/m< S^2$. Every arm can therefore independently become optimal, requiring its classical information allocation. For the path, use \eqref{alignment_paper:eq:jensenF} with $s=1$ and $\rho_L=1/\sqrt2$.
\end{proof}
This separation concerns the graph-energy classes, not superiority over all existing spectral algorithms. A bound based solely on $S\sqrt D$ gives the same width $\sqrt2\Gamma$ for both graphs and cannot express it. Nor do the normalized paths and cliques have identical Laplacian spectra: the example isolates the information lost by diameter compression.

\begin{figure}[t]
\centering
\begin{tikzpicture}
\begin{axis}[width=.9\linewidth,height=6cm,ymode=log,
 xlabel={Normalized supplied energy radius $S/\Gamma$},
 ylabel={Information constant divided by $2\sigma^2/\Gamma$},
 grid=major,legend columns=2,legend style={font=\small,at={(0.5,1.03)},anchor=south},xmin=0,xmax=1.3]
\addplot[red,thick] table[x=s,y=clique_constant_ratio,col sep=comma]{evidence/alignment_paper/results/geometry/information.csv};
\addlegendentry{Clique: exact information value}
\addplot[blue,thick,dashed] table[x=s,y=path_constant_upper,col sep=comma]{evidence/alignment_paper/results/geometry/information.csv};
\addlegendentry{Path: information upper bound}
\addplot[black,dotted,thick] table[x=s,y=classical_ratio,col sep=comma]{evidence/alignment_paper/results/geometry/information.csv};
\addlegendentry{Unstructured value}
\end{axis}
\end{tikzpicture}
\caption{Analytical information comparison at $m=256$. The path curve is $\min\{m,(1-s/\sqrt2)^{-2}\}$ and bounds the information-program value from above; it is not an attained policy coefficient or an exact path expression. At $s=1$ the clique value is $256$, while the path value is at most $11.66$.}
\label{alignment_paper:fig:geometryinformation}
\end{figure}



\subsubsection{A closed-form answer for homogeneous cliques}
The general program is difficult. The following solvable graph family gives a quantitative answer to the motivating question. The optimal component is a singleton. Each other component $c$ is a clique of $m_c\ge2$ arms, with common edge weight $w_c>0$. At the true instance its arms all have mean $a_c<\mu^\star$, so $\Gamma_c=\mu^\star-a_c$. Alternatives satisfy the energy certificate $S_c$ but need not be homogeneous. Define
\begin{equation}
 q_c=\frac{S_c}{\Gamma_c\sqrt{w_c(m_c-1)}},\qquad
 g_c=1+(m_c-1)\pos{1-q_c}^2.
 \label{alignment_paper:eq:alignmentratio}
\end{equation}

\begin{theorem}[Clique information constant]
For the graph-energy class just described, the value of \eqref{alignment_paper:eq:info} is
\begin{equation}
 C_{L,S}(\mu)=\sum_{c\ne c^\star}
 \frac{2\sigma^2 m_c\Gamma_c}
 {\Gamma_c^2+(m_c-1)
       \pos{\Gamma_c-S_c/\sqrt{w_c(m_c-1)}}^2}
 =\sum_{c\ne c^\star}\frac{2\sigma^2m_c}{\Gamma_c g_c}.
 \label{alignment_paper:eq:cliqueconstant}
\end{equation}
Thus $g_c$ is the reduction factor relative to that component's unstructured Gaussian information constant $2\sigma^2m_c/\Gamma_c$.
\label{alignment_paper:thm:clique}
\end{theorem}
\begin{proof}
Constraints separate by suboptimal component: alternatives raising a single component must be rejected, and any alternative raising multiple components contains at least one such information requirement. The true means, graph, and regret coefficients in a component are invariant to permutations of its arms. Averaging any feasible allocation over permutations preserves feasibility and cost, because each information constraint is linear in the allocation and the full alternative set is invariant. Hence an optimal component allocation can be uniform, $x_i=X/m$.

Fix an arm that an alternative raises above $\mu^\star$. At the infimum it is raised to $\mu^\star$, a movement $\Gamma$. The boundary is sufficient: scaling a nonnegative displacement toward the true constant vector reduces both its information cost and its graph energy. Displacements can first be clipped to $[0,\Gamma]$, which does not increase energy or squared distance. By convexity and symmetry, the other $m-1$ movements may be equal to $h$. The energy constraint is then
\begin{equation}
 w(m-1)(\Gamma-h)^2\le S^2,
\end{equation}
and the smallest squared distance occurs at
$h=\pos{\Gamma-S/\sqrt{w(m-1)}}$.
The infimal information per unit total allocation is
$[\Gamma^2+(m-1)h^2]/(2\sigma^2m)$.
Strictly improving alternatives approach this boundary continuously; when $S=0$, raise every arm in the component by the same amount instead. Since $\mu^\star<1$, the approximating alternatives remain in the mean domain. The smallest feasible $X$ is the reciprocal of this information rate. Multiplying by the common regret coefficient $\Gamma$ proves the result.
\end{proof}

At $q_c=0$, $g_c=m_c$ and the information cost becomes $2\sigma^2/\Gamma_c$: the clique behaves like one arm. For $q_c>1$, an individual arm can plausibly exceed the optimum while its peers retain their true means. At $q_c=1$ the same information constant follows by a limiting alternative with arbitrarily small peer changes. Thus $g_c=1$ for $q_c\ge1$, and the unstructured constant returns. The transition is continuous and depends on uncertainty relative to the reward gap.

For $m=8$, values $q=0,0.25,0.5,0.75,1$ give factors $8,4.9375,2.75,1.4375,1$. These are ratios of information lower-bound constants. They are not asserted finite-horizon speedups of SP-UCB.

The actual graph energy at the homogeneous true instance is zero even when a positive $S_c$ is supplied. In this theorem, $S_c$ measures how much mismatch the learner must allow, not the observed graph error. A true graph and a graph with a trustworthy tight certificate convey different amounts of information than the same graph with no certified relationship to rewards.

For a clique, Lemma~\ref{alignment_paper:lem:resistance} gives
\begin{equation}
 \varepsilon_c\le S_c\sqrt{\frac{2}{w_cm_c}}
 =\Gamma_c q_c\sqrt{\frac{2(m_c-1)}{m_c}},
 \label{alignment_paper:eq:cliquewidth}
\end{equation}
before clipping at one. Combining this with Theorem~\ref{alignment_paper:thm:upper} yields an explicit attainable upper bound in a strong-alignment regime. The information expression remains finite near $q_c=1$, whereas the pooled-only bound can become loose; its intersection with $A_c$ retains the classical guarantee. The different constants and thresholds expose the gap between the simple policy and an information-optimal allocation.



\subsubsection{The exact-alignment endpoint is achievable}
If every $S_c=0$, connectedness forces a constant mean on each component. Selecting any representative reduces the problem to an ordinary $M$-armed bandit. In the Gaussian case with a unique best component, a classical asymptotically optimal algorithm on these representatives achieves
\begin{equation}
 R_T=\left(\sum_{c\ne c^\star}\frac{2\sigma^2}{\Gamma_c}\right)\log T
      +o(\log T).
 \label{alignment_paper:eq:exact}
\end{equation}
The corresponding change-of-measure lower bound shifts all arms in a component together, and matches this expression. The reduction is thus attainable, independently of SP-UCB's conservative confidence constants. For bounded reward classes with $2\le M\le T$, the exact-component minimax rate is $\Theta(\sqrt{MT})$: a representative-based minimax policy supplies the upper bound, and a Bernoulli subclass with arbitrary component means supplies the lower bound~\cite{collected20,collected22}. For $M=1$, regret is zero.

This endpoint concerns exact equality of the means. Approximate similarity leaves a separate task of identifying the best member of a competitive component. The $A_c$ term for optimal components in Theorem~\ref{alignment_paper:thm:upper} makes that cost explicit.


\subsection{A policy that retains resistance geometry}
The information separation above motivates a policy that uses a graph-dependent design rather than a component-wide diameter. The goal here is an implementable finite-time result and an achieved separation in arm-count dependence. We do not claim a matching leading information constant for general graphs.

\subsubsection{Design-based component bounds}
Fix $p_c\in\mathscr P_{|C_c|}$ before observing rewards and compute
\begin{equation}
 a_c=S_c\kappa_{L_c}(p_c),\quad
 J_c=\{i\in C_c:p_{c,i}>0\},\quad d_c=|J_c|.
 \label{alignment_paper:eq:designbias}
\end{equation}
For a singleton set $a_c=0$ and $p_c=1$. These are known policy inputs. The graph-energy bound implies
\begin{equation}
 |\mu_i-p_c^\top\mu_c|\le a_c\quad(i\in C_c),\qquad
 p_c^\top\mu_c\le v_c\le p_c^\top\mu_c+a_c.
 \label{alignment_paper:eq:designtransfer}
\end{equation}
The important point is that a convex combination of observations can transfer information more accurately than a bound on every pair of arms separately. A numerical design can use the upper radius in Lemma~\ref{alignment_paper:lem:designradius}, preserving validity without exact optimization.

For a design target $n\ge1$, collect
$n_{c,i}(n)=\lceil np_{c,i}\rceil$ observations from each $i\in J_c$, and let
\begin{equation}
 \widehat m_c(n)=\sum_{i\in J_c}p_{c,i}\bar Y_i(n_{c,i}(n)),\qquad
 h_c(n)=\sigma\sqrt{2\ell\sum_{i\in J_c}
                     \frac{p_{c,i}^2}{n_{c,i}(n)}}\le b(n).
 \label{alignment_paper:eq:designestimate}
\end{equation}
The integer rounding costs at most $d_c$ additional observations. For this policy we assume independent iid arm noise streams with sub-Gaussian proxy $\sigma$, an explicit strengthening of the conditional-noise model sufficient for the weighted fixed-count averages. Component scheduling and elimination may still be adaptive. The target counts themselves are deterministic given the stage and fixed design.

\begin{lemma}[Design confidence and component optimism]
With probability at least $1-\delta$, simultaneously over all scheduled targets and ordinary arm counts through $T$,
\begin{equation}
 |\widehat m_c(n)-p_c^\top\mu_c|\le h_c(n),\qquad
 |\bar Y_i-\mu_i|\le b(n_i).
 \label{alignment_paper:eq:designevent}
\end{equation}
On this event, $\widehat m_c-h_c$ is a lower bound on $\mu^\star$, and $\widehat m_c+h_c+a_c$ is an upper bound on $v_c$.
\label{alignment_paper:lem:designconfidence}
\end{lemma}
\begin{proof}
For any fixed target the weighted noise average is sub-Gaussian with variance proxy $\sigma^2\sum_i p_{c,i}^2/n_{c,i}$. Its two-tail failure probability at the stated radius is at most $2e^{-\ell}$. There are at most $T$ possible targets per component and $T$ counts per arm. The union bound with \eqref{alignment_paper:eq:radius} gives the event. This can be defined on the independent potential arm streams even for targets an adaptive policy never collects; conditioning on collection is unnecessary. The final assertions follow from \eqref{alignment_paper:eq:designtransfer}.
\end{proof}

\subsubsection{Capped graph-design elimination}
Call the policy graph-design elimination with UCB fallback (GDE-UCB). Let $a_{\max}=\max_ca_c$. If $a_{\max}=0$, every component has a constant mean and we use ordinary UCB on one representative per component. Otherwise set
\begin{equation}
 n_{\rm cap}=\min\left\{T,\ 1+
          \left\lceil\frac{512\sigma^2\ell}{a_{\max}^2}\right\rceil\right\}.
 \label{alignment_paper:eq:designcap}
\end{equation}
The numerical constant is conservative and is fixed before running the experiments; it is not tuned using the true gaps.
\begin{enumerate}
 \item Start with all components active and target $n=1$.
 \item Bring every active component to the counts in \eqref{alignment_paper:eq:designestimate}. If completing this stage would exceed the remaining horizon, end the design phase before starting it.
 \item Compute $L_c=\widehat m_c-h_c$ and $U_c=\widehat m_c+h_c+a_c$. Remove components with $U_c<\max_{q\text{ active}}L_q$.
 \item If one component remains, or $n=n_{\rm cap}$, end the design phase. Otherwise replace $n$ by $\min\{2n,n_{\rm cap}\}$ and repeat.
 \item Run ordinary arm UCB on the remaining components, reusing all observations. If the only remaining component has $a_c=0$, any representative is optimal and can be selected throughout.
\end{enumerate}
The algorithm uses the graph, $S_c$, $\sigma$, $T$, and $\delta$. It does not use true means or gaps. The cap prevents endless design sampling inside a heterogeneous optimal component. Fallback preserves exploration of unsampled members in a surviving component.

\begin{theorem}[Finite-time regret with a resistance design]
Fix a component $c^\star$ containing an optimal arm and put
\begin{equation}
 g_c=\Gamma_c-a_c-a_{c^\star},\qquad
 d_c^{\rm reg}=\mu^\star-p_c^\top\mu_c,
 \quad Q_c=\sum_{i\in J_c}\Delta_i.
\end{equation}
When $a_{\max}>0$, define
\begin{equation}
 B_c=\begin{cases}
 \min\{n_{\rm cap},\ 1+64\sigma^2\ell/g_c^2\},&g_c>0,\\
 n_{\rm cap},&g_c\le0.
 \end{cases}
 \label{alignment_paper:eq:designBc}
\end{equation}
On the event \eqref{alignment_paper:eq:designevent}, no optimal component is eliminated, design-phase regret from $c$ is at most $B_cd_c^{\rm reg}+Q_c$, and total regret obeys
\begin{equation}
 \mathcal R_T\le\sum_c(B_cd_c^{\rm reg}+Q_c)
       +\sum_{c\in\mathcal A_{\rm end}}A_c.
 \label{alignment_paper:eq:designregret}
\end{equation}
Here $A_c$ is \eqref{alignment_paper:eq:Ac} and $\mathcal A_{\rm end}$ is the set of surviving components. For an expected-regret bound one may replace the last sum by $\sum_cA_c$ and add $\delta T$.
\label{alignment_paper:thm:designregret}
\end{theorem}
\begin{proof}
All lower bounds are at most $\mu^\star$ and an optimal component has upper bound at least $\mu^\star$, so elimination cannot remove it. At a completed target $n$,
$U_c\le v_c+2b(n)+a_c$ and
$L_{c^\star}\ge\mu^\star-a_{c^\star}-2b(n)$.
Thus $g_c>4b(n)$ suffices for elimination, which holds for
$n>32\sigma^2\ell/g_c^2$. The doubling schedule overshoots this threshold by at most a factor two, with the first-stage and cap corrections represented by \eqref{alignment_paper:eq:designBc}. If the horizon prevents the next stage, previously completed counts are smaller and obey the same bound.

At a target at most $B_c$, component regret is bounded by
$\sum_{i\in J_c}(B_cp_{c,i}+1)\Delta_i=B_cd_c^{\rm reg}+Q_c$.
During fallback an optimal arm remains available. On ordinary confidence, a selected suboptimal arm satisfies $\Delta_i\le2b(n_i)$ before its pull. Its additional pulls are therefore at most $1+8\sigma^2\ell/\Delta_i^2$, giving $A_c$ as a conservative bound even when pilot observations are reused. Add the two phases. Outside the event regret is at most $T$.
\end{proof}
For any fixed valid instance with $a_{\max}>0$, $\delta=1/T$ gives $n_{\rm cap}=O(\log T)$ and expected regret $O(\log T)$. The all-zero case reduces to the classical component problem. This establishes uniform efficiency on the supplied class, not a minimax bound or an optimal logarithmic coefficient. The dependence on $a_{\max}$ can be conservative when components have very different certificates.

\subsubsection{An achieved topology separation}
\begin{corollary}[A path can be rejected using two arms]
Consider the normalized path of Corollary~\ref{alignment_paper:cor:informationseparation}, one optimal singleton, homogeneous path mean $a$, and supplied $S=\Gamma=\mu^\star-a>0$. Use the endpoint design. If $T\ge2n_{\rm cap}+3$, GDE-UCB eliminates the path by its capped stage on \eqref{alignment_paper:eq:designevent}, never samples its interior, and satisfies
\begin{equation}
 R_T\le\frac{1024\sigma^2\ell}{\Gamma}+4\Gamma+\delta T.
 \label{alignment_paper:eq:achievedpath}
\end{equation}
For the clique with the same $m$, means, $S$, and resistance diameter, every uniformly efficient policy on its graph-energy class has
$\liminf R_T/\log T\ge2\sigma^2m/\Gamma$.
\label{alignment_paper:cor:achievedseparation}
\end{corollary}
\begin{proof}
For the path $a_c=\Gamma/\sqrt2$ and the optimal singleton has zero bias. At the uncensored cap, $b(n_{\rm cap})<a_c/16$. Hence
$a_c+4b(n_{\rm cap})<5\Gamma/(4\sqrt2)<\Gamma$, which forces elimination. The budget condition permits the stage to finish and prevents the cap from being censored by $T$. Path observation counts are at most $n_{\rm cap}+2$ and every such pull costs $\Gamma$; singleton observations cost zero. Since $n_{\rm cap}\le2+1024\sigma^2\ell/\Gamma^2$, the displayed bound follows. The clique assertion is Corollary~\ref{alignment_paper:cor:informationseparation} and \eqref{alignment_paper:eq:lower}.
\end{proof}
The path upper bound depends on $m$ through the logarithm in $\ell$, while the clique lower-bound coefficient grows linearly in $m$. For fixed $m$ and increasing horizon this gives an achieved separation in the arm-count dependence of the logarithmic term. The constants are loose, and this result does not assert that GDE-UCB reaches the path information upper bound in \eqref{alignment_paper:eq:pathinformation}. It also does not compare different policies on the same unrestricted reward class: the two graphs impose different energy constraints.

\subsubsection{The limit of treating a graph as untrusted advice}
\begin{proposition}[No free recovery of the structured logarithmic constant]
Fix an arbitrary candidate graph. For Gaussian rewards with variance $\sigma^2$, a policy uniformly efficient over all of $[0,1]^K$ satisfies, at every instance with a unique best arm and $\mu^\star<1$,
\begin{equation}
 \liminf_{T\to\infty}\frac{R_T}{\log T}
       \ge\sum_{i:\Delta_i>0}\frac{2\sigma^2}{\Delta_i},
 \label{alignment_paper:eq:untrustedlower}
\end{equation}
even if the true means are exactly constant on each candidate component.
\label{alignment_paper:prop:untrusted}
\end{proposition}
\begin{proof}
For any suboptimal $i$, the unrestricted class includes an alternative changing only $\mu_i$ to $\mu^\star+\eta<1$. The history change-of-measure constraint gives
$\liminf\E[n_i(T)]/\log T\ge2\sigma^2/(\Delta_i+\eta)^2$.
Let $\eta\downarrow0$, multiply by $\Delta_i$, and sum. Other arm observations cannot distinguish this alternative.
\end{proof}
This is a direct application of the established lower-bound framework, not a new impossibility principle~\cite{collected11,collected22}. Clus-UCB also discusses the corresponding obstruction to learning unrestricted widths~\cite{collected19}. It excludes a promise of the smaller structured logarithmic constant together with logarithmic regret on every arbitrary graph violation using only standard bandit feedback. It does not exclude finite-horizon benefits, weaker robustness guarantees, restricted graph errors, or additional informative data. The perturbation certificate in Proposition~\ref{alignment_paper:prop:perturb} is one explicit way to restrict such errors.



\subsection{Constructing a graph with a valid certificate}
There are three distinct quantities: the constructed graph, its actual agreement with unknown rewards, and the certificate supplied to the learner. Our guarantees concern their justified relationship, not a retrospective oracle score alone.

\paragraph{Known reward continuity.}
Suppose arm features $x_i$ have a known metric $d$ and reward continuity bound
$|\mu_i-\mu_j|\le H d(x_i,x_j)$ with known $H$. Any proposed component has the valid certificate
\begin{equation}
 S_c^2=H^2\sum_{i<j\in C_c}w_{ij}d(x_i,x_j)^2.
\end{equation}
A direct diameter certificate $\varepsilon_c\le H\max_{i,j\in C_c}d(x_i,x_j)$ is also available. Graph construction should balance having few components against creating components with a diameter small relative to their reward gaps. Since the gaps are unknown, using them in the construction would require a separate adaptive analysis.

\paragraph{Pilot observations.}
Uniform raw confidence intervals give
$|\mu_i-\mu_j|\le|\bar Y_i-\bar Y_j|+b(n_i)+b(n_j)$.
Thus one can obtain simultaneous edge-difference or diameter bounds, build a partition from a pilot phase, and freeze it. Restarting the main policy on fresh observations permits conditioning on that fixed partition, adding the pilot regret and its certificate failure probability. This does not assume that data-dependent pooled histories are automatically unbiased. A pilot of $n_0$ pulls per arm can cost as much as $Kn_0$ regret, potentially consuming the savings that pooling seeks to provide.

\paragraph{Graphs updated online.}
Our theorem does not cover a graph or partition that changes during exploitation. The confidence event for a fixed component does not automatically cover arbitrary subsets selected from the same reward history. A finite prespecified family can be covered simultaneously, with an appropriate union penalty; a general learned graph needs a dedicated adaptive guarantee. Including raw bounds in a minimum does not remove this issue.

\paragraph{Ordinal side information.}
If prior information only guarantees a rank-preserving filter, it does not imply \eqref{alignment_paper:eq:energy} with a small known $S_c$. An ordinal model and a graph-energy model are distinct restricted classes. Translating a particular aligned-graph definition into a regret reduction requires identifying what alternatives it excludes and how much information remains necessary to reject them.


\subsection{Experiments isolating the value of graph geometry}
\subsubsection{Protocol and model-matched width comparators}
The main experiment directly tests the path--clique construction. Each instance has one singleton of mean $0.8$ and one component of $m\in\{16,64,256\}$ arms with homogeneous mean $0.6$. The learner receives the fixed prior energy radius $S=0.2$ for that component, rather than its zero true energy. The algorithm does not receive the gap. Path edge weights are $(m-1)/2$ and clique weights are $1/m$. Both have resistance diameter two and the same implied diameter certificate $\varepsilon=0.2\sqrt2$. The graph and certificate are supplied, not estimated from selected-arm feedback.

There are 40 independent replications at horizon $T=20{,}000$, with Gaussian noise $\sigma=0.05$ and $\delta=1/T$. Policies use the same $\ell$ from \eqref{alignment_paper:eq:radius}, and share independent per-arm noise streams by arm identity and pull number within a replication. The implementation is standard-library Python and stores per-run counts, pilot counts, eliminated components, confidence diagnostics, and trajectories. Cross-policy and cross-graph results within a seed are coupled; independence is across replications within a condition. Error bars are $1.96$ standard errors of simulated mean regret.

We compare ordinary UCB, the diameter-based SP-UCB policy, GDE-UCB, and a Gaussian width-profile UCB motivated by Clus-UCB~\cite{collected19}. The last comparator uses the index
\begin{equation}
 u_i^{\rm width}=\sup\left\{q:
 n_i(q-\bar Y_i)_+^2+
 \sum_{\substack{j\in C_c\\j\ne i}}n_j(q-\bar Y_j-\varepsilon_c)_+^2
 \le2\sigma^2\ell\right\},\quad i\in C_c.
 \label{alignment_paper:eq:widthprofile}
\end{equation}
It initializes every arm. This is an explicitly specified one-sided Gaussian analogue, not a reproduction of the published Bernoulli algorithm or a claim that its asymptotic theorem transfers unchanged. The upper root is solved as a piecewise quadratic using sorted breakpoints, with an independent bisection check. Both width-only methods receive exactly the same inputs for the two normalized graphs, so their trajectories are identical under the shared streams. These identical conditions are replayed in the output rather than treated as independent datasets.

GDE-UCB uses the analytic endpoint design for paths and uniform design for cliques, with cap constant 512 as in \eqref{alignment_paper:eq:designcap}. No policy parameters are fitted to realized rewards or selected retrospectively. The design phase is exploratory, so its cost is expected to outweigh its benefit on some small instances.

\begin{table}[t]
\centering\small
\begin{tabular}{lrrrrr}
\toprule
Graph & $m$ & UCB & Width profile & GDE-UCB & GDE arms\\
\midrule

Path & 16 & $10.81\pm0.21$ & $10.80\pm0.21$ & $40.32\pm3.97$ & 2 \\
Path & 64 & $45.69\pm0.37$ & $45.38\pm0.44$ & $40.32\pm3.97$ & 2 \\
Path & 256 & $191.83\pm0.96$ & $187.15\pm2.26$ & $46.08\pm3.21$ & 2 \\
Clique & 16 & $10.81\pm0.21$ & $10.80\pm0.21$ & $163.20\pm0.00$ & 16 \\
Clique & 64 & $45.69\pm0.37$ & $45.38\pm0.44$ & $166.40\pm0.00$ & 64 \\
Clique & 256 & $191.83\pm0.96$ & $187.15\pm2.26$ & $216.64\pm0.51$ & 256 \\


\bottomrule
\end{tabular}
\caption{Mean final regret with $1.96$ standard-error half-widths; the final column is the mean number of distinct sampled suboptimal arms. SP-UCB's values equal UCB's in these conditions and are retained in the CSV data. The graph-energy certificate is the same positive $S=0.2$ in every main condition.}
\label{alignment_paper:tab:geometryresults}
\end{table}

\begin{figure}[t]
\centering
\begin{tikzpicture}
\begin{axis}[width=.94\linewidth,height=7cm,xmode=log,log basis x=2,
 xlabel={Suboptimal component size $m$},ylabel={Mean final cumulative regret},
 grid=major,legend columns=2,legend style={font=\small,at={(0.5,1.03)},anchor=south},xtick={16,64,256}]
\addplot[black,thick,mark=*,error bars/.cd,y dir=both,y explicit]
 table[x=m,y=ucb,y error=ucb_ci,col sep=comma]{evidence/alignment_paper/results/geometry/path_scaling.csv};
\addlegendentry{UCB, either graph}
\addplot[teal,dashed,thick,mark=triangle*,error bars/.cd,y dir=both,y explicit]
 table[x=m,y=width_profile,y error=width_profile_ci,col sep=comma]{evidence/alignment_paper/results/geometry/path_scaling.csv};
\addlegendentry{Gaussian width profile, either graph}
\addplot[blue,thick,mark=square*,error bars/.cd,y dir=both,y explicit]
 table[x=m,y=design,y error=design_ci,col sep=comma]{evidence/alignment_paper/results/geometry/path_scaling.csv};
\addlegendentry{GDE-UCB, path}
\addplot[red,thick,mark=diamond*,error bars/.cd,y dir=both,y explicit]
 table[x=m,y=design,y error=design_ci,col sep=comma]{evidence/alignment_paper/results/geometry/clique_scaling.csv};
\addlegendentry{GDE-UCB, clique}
\end{axis}
\end{tikzpicture}
\caption{The graph geometry matters even when diameter certificates are identical. GDE-UCB samples only the two path endpoints, whereas its uniform clique design incurs substantial overhead. The figure supports the achieved path regime, not universal superiority of the policy.}
\label{alignment_paper:fig:geometryscaling}
\end{figure}



\subsubsection{Results, overhead, and certificate failure}

On paths, GDE-UCB's mean final regret is 40.32 at $m=16$ and 46.08 at $m=256$, compared with ordinary UCB's 10.81 and 191.83, respectively. At the larger size the observed UCB/GDE-UCB regret ratio is 4.16; at the smaller size the design overhead makes GDE-UCB worse. This finite-horizon comparison supports the predicted arm-count dependence, not a universal speedup or an optimality claim.

GDE-UCB samples exactly two suboptimal path arms in every condition and eliminates the path in every replication. The other policies sample every suboptimal arm. On cliques, the uniform design does not eliminate the component in these runs; it samples all its arms and incurs mean regret 216.64 at $m=256$. Thus the same policy also exposes the cost of an unhelpful design regime. The width-profile comparator remains close to ordinary UCB because the shared width certificate cannot exclude independent near-optimal alternatives.



These comparisons isolate a failure of diameter compression. They do not imply that the graph-design algorithm dominates spectral bandits. In particular, the spectral baseline in the supplementary experiment uses a different exact certificate, and is not a matched comparison in this positive-energy experiment. A tuned spectral comparison on the same energy class, more varied graphs, and application datasets remain needed for an empirical superiority claim.

A separate stress test makes the assumption boundary explicit. The graph is a path of eight arms with means $(0.4,0.4,0.4,0.9,0.4,0.4,0.4,0.4)$ plus a singleton of mean $0.8$. A valid supplied radius is computed from this synthetic vector's energy and labeled oracle information. An invalid zero radius forces the representative-based policy to ignore the hidden optimal arm. Table~\ref{alignment_paper:tab:certificatestress} compares two horizons; its approximately constant regret per round in the invalid case demonstrates linear loss. Ordinary observation confidence can hold throughout while the structural assumption fails.

\begin{table}[t]
\centering\small
\begin{tabular}{lrrr}
\toprule
Certificate & Horizon & Mean regret & Regret per round\\
\midrule

Valid oracle & 10000 & 22.80 & 0.0023 \\
Invalid zero & 10000 & 1000.42 & 0.1000 \\
Valid oracle & 20000 & 24.00 & 0.0012 \\
Invalid zero & 20000 & 2000.43 & 0.1000 \\


\bottomrule
\end{tabular}
\caption{Certificate stress test. The valid certificate is explicitly oracle-supplied; the invalid certificate asserts exact equality despite a hidden peak. This experiment illustrates a limitation, not a method for learning the correct certificate.}
\label{alignment_paper:tab:certificatestress}
\end{table}

Numerical checks cover the resolvent's least informative alternatives, clique agreement, resistance-radius primal/dual gaps, the first-order inequalities on a nonsymmetric weighted graph, confidence coverage, and fallback when the optimal component is heterogeneous. They also compare the width-profile oracle to direct scalar bisection. These checks help detect implementation and algebra errors; they do not replace independent proof review.



\subsection{Discussion and remaining limits}
The graph's value is the information it excludes about competing reward vectors. A diameter certificate provides a simple sufficient condition for sharing observations, but it is not a sufficient statistic for that information. The path--clique separation demonstrates a factor growing with component size at identical diameter and supplied energy. The resistance radius provides an interpretable local sensitivity parameter, while the full information design captures behavior beyond the local regime.

The algorithmic separation is attained under quantitative side information. Endpoint sampling rejects a homogeneous normalized path without exploring its interior, and ordinary UCB fallback handles a heterogeneous surviving component. This gives an implementable logarithmic guarantee with conservative constants. It does not give an information-optimal policy for every graph, a minimax-optimal bound, or empirical superiority over all spectral methods. The capped design can be expensive on small instances and on graphs requiring broad support.

Graph correctness has two distinct roles. A valid certificate determines the model class; its precision determines how much information transfer is useful. A certified Laplacian perturbation can be charged through Proposition~\ref{alignment_paper:prop:perturb}. A false certificate can invalidate optimism and yield linear regret, as the stress experiment shows. The unrestricted lower bound in Proposition~\ref{alignment_paper:prop:untrusted} explains why learning an arbitrary graph's validity cannot preserve the smaller structured logarithmic constant for free.

The most consequential remaining questions are precise. One is a computationally tractable policy attaining the general program \eqref{alignment_paper:eq:info} for nonhomogeneous graph-energy instances, with the regularity and boundary conditions needed by structured-bandit allocation tracking. Another is a finite-strength spectral-index regret theorem that accounts for adaptive counts and comparison bias. A third is a restricted model of uncertain graph information with a proved price for validation or repair. The present resolvent connection and certified perturbation comparison supply tools for these questions, not their complete solution.

The evidence supports a graph-specific theoretical contribution and a reproducible demonstration of its regime. It does not establish priority for each specialized theorem. Closest comparisons include spectral bandits, clustered bandits, GRUB's pure-exploration geometry, and misspecified kernel and linear models. A submission should be evaluated against those alternatives, with application data and tuned graph-aware baselines supplementing the controlled experiments.


\subsection{Supplementary certificate and graph-corruption experiments}
\subsubsection{Protocol}
The accompanying standard-library Python implementation runs 40 independent macro-replications with horizon $T=20{,}000$ and Gaussian observation noise $\sigma=0.15$. There are 25 arms: one with mean $0.8$, and three groups of eight arms with means $0.65$, $0.5$, and $0.35$. The optimal arm is a singleton. Clique edge weights are $1/m$, so each nonzero component Laplacian eigenvalue equals one. Policies share per-arm noise streams indexed by pull number within a replication; this coupling does not give any policy extra observations. Different macro-replications use independent seeds.

Ordinary UCB and SP-UCB use \eqref{alignment_paper:eq:radius}, with $\delta=1/T$. The regularized spectral baseline is described below. Results report realized pseudo-regret, averaged over replications. The interval half-width is $1.96$ times the across-run standard error; it is an approximate normal confidence interval for a simulation mean. Raw observations are Gaussian, while the regret variable remains bounded by $T$ because the means lie in $[0,1]$.

\subsubsection{Certificate sensitivity with fixed graph and rewards}
The first experiment leaves the graph and true rewards fixed and supplies SP-UCB with diameter bounds $\varepsilon=0,0.05,0.1,0.2,0.4$ on each suboptimal component. All are valid because those true components are homogeneous. This intervention measures the value of more precise side information, not a change in actual graph alignment or a learned certificate.

The graph-regularized baseline uses
\begin{equation}
 V_t=\diag(n_i(t))+V_0,\quad V_0=\eta I+\lambda L,\quad
 \widehat\mu_t=V_t^{-1}\sum_{s<t}e_{A_s}Y_s,
\end{equation}
and selects the largest
\begin{equation}
 \widehat\mu_{t,i}+\beta_t\sqrt{(V_t^{-1})_{ii}},\quad
 \beta_t=\sigma\sqrt{\log\frac{\det V_t}{\det V_0}+2\log\frac1\delta}
       +\sqrt{\eta K+\lambda\sum_cS_c^2}.
\end{equation}
This is a regularized SpectralUCB-style baseline, using the usual self-normalized linear confidence construction~\cite{collected17,collected23}. Here $\eta=0.01$, $\lambda=1000$, and the exact energy certificate is zero. It is a single fixed configuration, not a tuned reproduction of every published spectral algorithm. The prior-norm term is valid because $\|\mu\|^2\le K$ and the supplied energies are valid. This baseline already exploits graph information, so outperforming ordinary UCB alone would not establish superiority over existing graph methods.

\subsubsection{Corrupted graph construction}
The second experiment exchanges 0, 1, 2, or 4 members of the highest and lowest suboptimal groups, keeping group sizes eight and the optimal singleton intact. It then builds cliques on the resulting memberships. The graph therefore increasingly joins arms with different means. Each policy receives an explicitly \emph{oracle-supplied} energy certificate computed from the true synthetic means, converted to a diameter bound by Lemma~\ref{alignment_paper:lem:resistance}. This separates the cost of incorrect edges from the separate problem of certifying their error. It gives no evidence that such a certificate can be inferred cheaply from selected-arm feedback.

\begin{table}[t]
\centering\small
\begin{tabular}{llrr}
\toprule
Experiment & Policy/condition & Mean regret & 95\% half-width\\
\midrule

Certificate & ucb & 100.02 & 1.78 \\
Certificate & spectral\_ucb & 28.65 & 0.97 \\
Certificate & pool\_eps\_0.0 & 12.22 & 0.67 \\
Certificate & pool\_eps\_0.05 & 21.07 & 1.15 \\
Certificate & pool\_eps\_0.1 & 54.25 & 3.80 \\
Certificate & pool\_eps\_0.2 & 83.08 & 2.26 \\
Certificate & pool\_eps\_0.4 & 100.02 & 1.78 \\
Perturbed graph & ucb & 100.02 & 1.78 \\
Perturbed graph & pool\_swaps\_0 & 12.22 & 0.67 \\
Perturbed graph & pool\_swaps\_1 & 75.37 & 1.74 \\
Perturbed graph & pool\_swaps\_2 & 75.37 & 1.74 \\
Perturbed graph & pool\_swaps\_4 & 75.37 & 1.74 \\


\bottomrule
\end{tabular}
\caption{Results at $T=20{,}000$ across 40 macro-replications. Widths in the certificate experiment are supplied diameter bounds. Swap counts in the graph experiment change component memberships; true energy is then supplied as an oracle certificate.}
\label{alignment_paper:tab:results}
\end{table}

\begin{figure}[t]
\centering
\begin{tikzpicture}
\begin{axis}[width=.93\linewidth,height=7cm,xmode=log,
 xlabel={Horizon},ylabel={Mean cumulative pseudo-regret},grid=major,
 legend pos=north west,legend style={font=\small}]
\addplot[black,thick,mark=*,error bars/.cd,y dir=both,y explicit] table[x=t,y=ucb,y error expr={1.96*\thisrow{ucb_se}},col sep=comma]{evidence/alignment_paper/results/certificate_curves.csv};
\addlegendentry{Ordinary UCB}
\addplot[teal,thick,mark=triangle*,error bars/.cd,y dir=both,y explicit] table[x=t,y=spectral_ucb,y error expr={1.96*\thisrow{spectral_ucb_se}},col sep=comma]{evidence/alignment_paper/results/certificate_curves.csv};
\addlegendentry{Spectral baseline, exact certificate}
\addplot[blue,thick,mark=square*,error bars/.cd,y dir=both,y explicit] table[x=t,y=pool_eps_0.0,y error expr={1.96*\thisrow{pool_eps_0.0_se}},col sep=comma]{evidence/alignment_paper/results/certificate_curves.csv};
\addlegendentry{SP-UCB, $\varepsilon=0$}
\addplot[red,thick,mark=diamond*,error bars/.cd,y dir=both,y explicit] table[x=t,y=pool_eps_0.1,y error expr={1.96*\thisrow{pool_eps_0.1_se}},col sep=comma]{evidence/alignment_paper/results/certificate_curves.csv};
\addlegendentry{SP-UCB, $\varepsilon=0.1$}
\addplot[orange,thick,mark=+,error bars/.cd,y dir=both,y explicit] table[x=t,y=pool_eps_0.2,y error expr={1.96*\thisrow{pool_eps_0.2_se}},col sep=comma]{evidence/alignment_paper/results/certificate_curves.csv};
\addlegendentry{SP-UCB, $\varepsilon=0.2$}
\end{axis}
\end{tikzpicture}
\caption{Certificate sensitivity with a fixed exactly homogeneous graph. Curves average 40 replications; error bars show $1.96$ standard errors. The $\varepsilon=0.4$ final value equals the ordinary-UCB value in this experiment and is omitted for clarity. Final uncertainty estimates appear in Table~\ref{alignment_paper:tab:results}; full per-checkpoint standard errors are provided in the CSV data.}
\label{alignment_paper:fig:empirical}
\end{figure}

\subsubsection{Findings and their interpretation}
With exact certificates, SP-UCB's average final regret is $12.22$, compared with $100.02$ for ordinary UCB and $28.65$ for the configured spectral baseline. The observed ordinary-UCB/SP-UCB ratio is approximately $8.18$, consistent with rejecting groups of eight equivalent suboptimal arms. It is a finite-horizon observation, not a theorem about an eightfold reduction for arbitrary graphs.

As the supplied width becomes looser, average SP-UCB regret increases to $21.07$, $54.25$, $83.08$, and $100.02$. This supports the qualitative certificate dependence in Theorem~\ref{alignment_paper:thm:upper}. It does not show that its constants are tight or that empirical regret is monotone in every instance.

Exchanging even one high/low member raises mean regret to $75.37$. Further exchanges give the same final summary in this construction. The supplied width of both affected components is then large enough that pooling offers no useful rejection bound, while the unchanged middle component still benefits. This plateau is informative: graph-error severity and regret need not have a smooth one-to-one relationship. Once the relevant alignment threshold is crossed, additional edge corruption may have little effect on this particular index.

The implementation checks ordinary-arm and pooled-noise confidence coverage, and verifies the arm and component pull-count implications on covered paths. Analytical checks confirm rank cancellation for uniform complete graphs, the clique least-informative alternative, endpoint constants, and the energy-to-width certificates for corrupted partitions. These checks support implementation consistency; they are not a substitute for mathematical peer review.

\subsubsection{Scope of the numerical evidence}
The experiments use a favorable family with an isolated optimal arm and homogeneous true suboptimal groups. They assess certificates and graph corruption in a controlled setting. They do not include online graph learning, near-tied heterogeneous optimal components, temporal dependence, real applications, or a full tuning study. Clus-UCB and hierarchical Thompson sampling are close competitors whose implementation and careful comparison would be needed before an empirical superiority claim. The manuscript makes no such claim.

\subsection{Reproducibility and numerical implementation}
The main geometry experiment is generated by
\texttt{experiments/resistance\_geometry.py}; supplementary experiments use
\texttt{experiments/alignment\_regret.py}. Both require only Python's standard library. From the project root:
\begin{verbatim}
python3 experiments/alignment_regret.py --verify
python3 experiments/alignment_regret.py --runs 40 --horizon 20000
python3 experiments/resistance_geometry.py --verify
python3 experiments/resistance_geometry.py --runs 40 --horizon 20000
python3 experiments/render_geometry_results.py
cd research/alignment_paper
tectonic --only-cached --keep-logs manuscript.tex
\end{verbatim}
The output directory includes per-replication regrets and counts, aggregate means and standard errors, curves at recorded checkpoints, the analytical lower-bound curve, and metadata specifying seeds, mean vectors, graph memberships, certificates, policy parameters, and a SHA-256 hash of the experiment source. The figures use PGFPlots directly on these CSV files. Results in the manuscript are actual generated values; no simulation tables are manually fabricated.

For the spectral baseline, a clique's regularized design matrix is
$V_c=\diag(d_i)-a\one\one^\top$, where $d_i=n_i+\eta+\lambda$ and $a=\lambda/m$ under edge weights $1/m$. Let $h_i=1/d_i$ and $b=a/(1-a\sum_i h_i)$. Then
\begin{equation}
 V_c^{-1}=\diag(h_i)+b hh^\top,\qquad
 \log\det V_c=\sum_i\log d_i+\log(1-a\sum_i h_i).
\end{equation}
These identities avoid an external numerical dependency and a dense matrix inverse at each step. The same expression covers a singleton, where the graph penalty cancels and $V_c=n_i+\eta$.

\subsection{A sharper distinction between uncertainty and actual error}
Let $S_c^{\rm true}=\sqrt{\mu_c^\top L_c\mu_c}$. A valid supplied bound must satisfy $S_c\ge S_c^{\rm true}$. The true quantity evaluates a graph retrospectively; the supplied bound determines the alternatives the policy must consider. Tightening $S_c$ without violating validity reduces an information class. Changing $L_c$ can improve or worsen both $S_c^{\rm true}$ and effective resistance. The comparison is not determined by one of these quantities alone.

For example, adding a strong edge between substantially different means increases graph energy. Increasing all weights, in contrast, increases energy and decreases resistance by inverse factors, without changing the certified width when the certificate is rescaled consistently. A graph-quality evaluation should therefore report component sizes, gap scales, true diagnostic energy, supplied certificate size, and the resulting reward-diameter bound. Only the latter information available before a decision can justify its confidence interval.

